\chapter{LANDASAN TEORI}

\vspace{0.5cm}
\noindent\fbox{\begin{minipage}{\textwidth}
\textit{\textbf{Pedoman Penulisan:} Bab landasan teori merupakan tinjauan pustaka yang menguraikan teori-teori serta algoritma yang mendukung judul. Landasan teori bukan hanya berupa definsi-definisi, melainkan konsep dasar tentang teori tersebut serta model atau algoritma yang berkaitan langsung dengan masalah yang diteliti. Pada bab ini juga ditulis tentang tools/software komponen yang digunakan untuk aplikasi atau untuk keperluan penelitian.}
\end{minipage}}
\vspace{0.5cm}


\section{Tinjauan Pustaka}

\vspace{0.5cm}
\noindent\fbox{\begin{minipage}{\textwidth}
\textit{\textbf{Pedoman Penulisan:} Berisikan penjelasan tentang konsep dasar program, metode algoritma, pengujian aplikasi dan peralatan pendukung yang digunakan dalam merancang aplikasi.\\
Catatan:\\
Peralatan pendukung (tools system) menggunakan Unified Modelling Languange (UML), menjelaskan teori UML dan diagram-diagram yang digunakan yaitu Activity, Use Case, Sequence, Class dan Deployment, tapi jika dalam pembuatan perangkat lunak tidak menggunakan class maka teori dan pembahasan hanya meliputi empat diagram UML saja.}
\end{minipage}}
\vspace{0.5cm}


\section{Penelitian Terkait}

\vspace{0.5cm}
\noindent\fbox{\begin{minipage}{\textwidth}
\textit{\textbf{Pedoman Penulisan:} a. Berisikan penjelasan tentang aplikasi yang akan dikembangkan berdasarkan penelitian serupa yang diambil dari jurnal-jurnal ilmiah.\\
b. Berisikan Kutipan dari jurnal yang membahas penelitian terkait dengan tema yang diambil. Menggunakan software Mendeley dengan metode APA. Jurnal terkait minimal 5 tahun kebelakang dari tahun berjalan. Minimal 5 kutipan jurnal terkait.}
\end{minipage}}
\vspace{0.5cm}

